Crypto-assets are issued, traded and described on the Web, and the vocabulary
that governs them is published as Web documents by regulators, standard-setters
and researchers who do not share a dictionary: the same token is a crypto-asset
under MiCA, a digital asset under US law and a virtual asset under FATF, and
each definition discriminates along dimensions the text never names. Manual
taxonomies capture one perspective and go stale, while LLM taxonomy and
ontology-learning methods assume a fixed label space and do not trace their
output to a legal source. In this paper we
propose \sys, a framework that induces the classification dimensions, which we
call term types, from the definitions themselves. \sys (1) builds a
definitions-first corpus from 24 legal instruments and academic taxonomies,
(2) extracts term types in a single pass over the whole corpus held in a cached
context, requiring a document, section and verbatim quote for every dimension
and value, (3) verifies every quote against the source, (4) orders the term
types into an is-a graph by asking, for each one, for a cited probability
distribution over the more general dimensions it specialises, consolidated
greedily into a directed acyclic graph, and (5) uses the graph as the label
space for Web content and grows it only where tagged content raises its
information gain. On a corpus spanning the EU, the UK, the
US and six international bodies, \sys produces 93 term types with 317 values and
480 citations, 97.9\% of which verify, records a cross-source disagreement for 69
of them, and links them by 279 weighted is-a edges under one root per category,
for a few dollars of API calls. Two maintenance iterations over \AllItems\ DLT-Corpus papers and patents
verify 96--97\% of tag evidence and grow the graph by \AllAddedTypes\
dimensions and \AllAddedValues\ values for \$\AllCost, and the dimensions
discovered first are among the most used next. \todo{One sentence of comparative evidence once RQ2--RQ4 are run.} The
corpus, taxonomy and code are available at an anonymised repository.\footnote{\todo{anonymised URL}}
